% Bilan d'ensemble (chiffres : sorties/architectures_ircsg.txt, sorties/perte_seche_autres.txt,
% ../foncier/sorties/simulation.txt, ../is/sorties/simulation_is.txt)

\begin{table}[!htbp]\centering\scriptsize
\caption{Les quatre réformes, à recette constante (Md€ par an, régime permanent)}\label{tab:ensemble}
\begin{tabularx}{\linewidth}{@{}>{\raggedright\arraybackslash}p{3.2cm}>{\raggedright\arraybackslash}p{2.6cm}>{\raggedright\arraybackslash}p{2.4cm}>{\raggedright\arraybackslash}X@{}}
\toprule
Réforme & Recette & Gain d'efficacité (fourchette) & Qui paie plus \\
\midrule
Fusion IR-CSG-CRDS : franchise d'activité 4~000~€, 6~000~€ par enfant, petites pensions exonérées & 0 par
construction ; moindre optimisation : +1,9 à +5,7 & 4,9 (2,3 à 11,1) & retraités (+15~Md€) ; 29~\% des foyers
plus de 500~€ \\
Immobilier : droits de mutation $-$50~\%, taxe foncière à deux taux, taux national sur le sol & 0 (transition :
hausse temporaire de la taxe actuelle) & 1,9 (1,0 à 4,2) & propriétaires de terrains chers (Paris,
Hauts-de-Seine) \\
IS : déduction immédiate de l'investissement, financée par l'assiette & 0 en régime permanent ; dette de
transition de 2 à 32 & 0,1 (0,05 à 0,16), moins 0 à 0,8 si fin du taux réduit & sociétés au taux réduit,
crédit d'impôt recherche \\
PFU maintenu, report d'apport-cession & 0 ; recette différée selon le traitement du report & 0 & selon le
traitement du report \\
\midrule
Ensemble & 0 & environ 7 (2,5 à 15,5) & \\
\bottomrule
\end{tabularx}
\end{table}

\textbf{Le gain total est modeste et robuste} (tableau~\ref{tab:ensemble}). Il est d'environ 7~Md€ par an au
centre, 0,2~\% du PIB. Il pourrait être trois fois moindre ou deux fois plus grand selon les comportements.
Dans tous les jeux d'hypothèses, la fusion et la réforme de l'immobilier réduisent la perte sèche, et aucun
gain ne repose sur une réaction déclarative. Plus des deux tiers du gain viennent de la fusion de l'impôt sur le revenu et de la CSG, un peu plus
d'un quart du remplacement des droits de mutation. L'IS et le PFU servent surtout à ne pas aggraver les choses :
ne pas monter des taux que les bénéfices et les dividendes savent fuir, et rendre l'impôt prévisible. Le gain
de la fusion représente environ deux tiers de ce qu'un barème optimisé sans contrainte obtiendrait. Le dernier
tiers demanderait de rendre plus de 20~Md€ au dernier décile, sur la foi d'élasticités que personne n'a
mesurées pour un tel changement.

\textbf{Les grandes fortunes.} Trois profils sont à distinguer.
\begin{itemize}\setlength{\itemsep}{0pt}
\item \emph{Le dirigeant salarié.} Son taux marginal passe de 55~\% environ (avec la CEHR) à 50~\%. Le 1~\% le
plus aisé paie 2,7~Md€ de moins sur ses revenus du barème.
\item \emph{Le rentier.} Il garde le PFU, majoration comprise. Il paie plus de taxe foncière s'il possède des
terrains chers, et moins de droits quand il achète.
\item \emph{L'entrepreneur qui a apporté ses titres à une holding.} Son sort dépend du traitement du report
d'apport-cession (section~\ref{sec:pfu}). L'expiration à toute transmission alourdirait de 90 à 100~M€ la
facture de chacun des plus gros apporteurs ; la transmission du report sans effacement n'augmente rien à la
transmission.
\end{itemize}
Une incertitude demeure pour les 3\,800 foyers au-delà de 3~M€ de revenu fiscal : selon la part (non publiée)
de leurs revenus au PFU, la fusion baisse leur impôt sur le barème ou l'augmente de plusieurs points. Il faut
la lever avec les données de la DGFiP avant d'arrêter le barème.

\textbf{Les familles} gagnent : 1\,200~€ par an en moyenne pour un couple avec deux enfants, 230~€ pour un
parent isolé. La franchise de 6\,000~€ par enfant vaut 6\,000~€ fois le taux marginal du parent : 840~€ par
enfant sous 20\,000~€ de revenu imposé, 1\,800~€ vers 45\,000~€, davantage au-delà, contre 742~€ en moyenne
pour le quotient familial actuel. Les droits de mutation, payés au moment où une famille achète plus grand,
sont divisés par deux.

\textbf{Les retraités} sont les principaux contributeurs. La fusion leur fait payer 15~Md€ de plus,
2\,300~€ par an pour un couple, car elle met fin à leurs taux réduits et à leur abattement. Propriétaires plus
souvent que la moyenne, ils paient aussi la hausse de taxe foncière, sans profiter de la baisse des droits
s'ils ne déménagent pas. C'est le prix du gain d'efficacité : les pensions ne réagissent pas à l'impôt, et les
taxer comme les salaires permet de baisser ce qui pèse sur le travail. Trois leviers atténuent cette charge :
\begin{itemize}\setlength{\itemsep}{0pt}
\item l'exonération des petites pensions, retenue ici ;
\item le maintien de l'abattement de 10~\%, pour 0,6~point de taux pour tous ;
\item une transition sur quatre ans, et, pour la taxe foncière, le report de paiement des propriétaires âgés.
\end{itemize}

\textbf{Les dates.} Les chiffres décrivent le droit de 2026. L'IR et la CSG sont mesurés sur les revenus de
2024, l'IS en euros de 2027, l'immobilier sur les ventes de 2023-2024 et la fiscalité locale de 2024.
